% Ad Atticum 11.15 --- headnote
% ad-atticum-11-15, 14 May 47 BC, Brundisium

\textit{Cicero to Atticus, written from Brundisium on
the day before the Ides of May 47 BC --- 14 May (the
manuscript dateline: \textit{Scr.\ Brundisi prid.\
Id.\ Mai.\ a.\ 707 (47)}). The colophon at the end
of the letter (\textit{Pr.\ Idus Maias}) confirms the
date. Atticus has explained, in good faith, why he
cannot come down to Brundisium just now; Cicero
accepts the reasons but presses the question that
they leave hanging --- what is he himself to do? The
geopolitical map he sketches in \S1 is the whole of
his predicament. Caesar is bogged down in the
Alexandrine war and apparently too embarrassed even
to write home about it; the Pompeians from Africa
look poised to come over; the men of the Achaian
party are about to leave Asia, either for Africa to
join them or for some neutral place. Cicero, with
perhaps one other man, has nowhere to go: no return
to the Pompeians, no real hope from the Caesarians.}

\textit{The middle of the letter is the most exposed
self-accusation in Book 11. Quintus has written
again, more bitterly than before, and Quintus's son
with an extraordinary hatred; every evil, Cicero
says, that can be invented is pressing him. But all
of it is easier to bear than the pain of having
done wrong --- of having returned to Italy and
accepted Caesar's protection --- which is at its
height (\textit{maximus}) and unending
(\textit{aeternus}). Every other man's case, he
argues, has an outlet: prisoners, the cut-off, even
those who went to Fufius (Caesar's lieutenant at
Patrae) of their own free will can at worst be
called timid; many more will be received back by
the Pompeians on any terms they like. Only his own
fault, and perhaps Laelius's, cannot be undone ---
and even the small comfort of Cassius's parallel
choice has just collapsed, since Cassius is now
said to have changed his plan of joining Caesar at
Alexandria. The closing two sections turn back to
the practical question --- should he creep nearer
to Rome in secret, or cross the sea? --- and to the
Fufidian inheritance, where Cicero suspects his
coheirs are stalling because they reckon his cause
is doomed. One short textual crux in \S3 is
preserved as a \dag\ marker; the sense given is the
most natural reading. ``Aesopus's son'' is the
disreputable boy of the famous tragic actor Clodius
Aesopus, of whom Cicero had reluctantly taken
charge.}
